\documentclass[10pt]{article} % For LaTeX2e
\usepackage{tmlr}
% If accepted, instead use the following line for the camera-ready submission:
%\usepackage[accepted]{tmlr}
% To de-anonymize and remove mentions to TMLR (for example for posting to preprint servers), instead use the following:
%\usepackage[preprint]{tmlr}

\usepackage{hyperref}
\usepackage{url}
\usepackage[utf8]{vietnam}
\usepackage{amsmath}
\usepackage{amssymb}
\usepackage{graphicx}
\usepackage{longtable}
\usepackage{multirow}
\usepackage{booktabs}
\usepackage{xcolor}
\usepackage[ruled,linesnumbered]{algorithm2e}
\SetKwInOut{KwRequire}{Require}
\SetKwInOut{KwEnsure}{Ensure}

% Bản nháp tiếng Việt: \tbd đánh dấu số liệu chờ thí nghiệm (paper/TMLR_pivot_plan.md),
% \cantrich đánh dấu chỗ cần bổ sung trích dẫn đã kiểm chứng DOI. Xoá cả hai trước khi nộp.
\newcommand{\tbd}[1]{\textcolor{red}{[TBD: #1]}}
\newcommand{\cantrich}[1]{\textcolor{orange}{[cần trích dẫn: #1]}}

\title{When Does Relational Distillation Transfer Across Electrode Montages? A Multi-Backbone Study for Low-Electrode SSVEP Recognition}

% Authors must not appear in the submitted version. They should be hidden
% as long as the tmlr package is used without the [accepted] or [preprint] options.
% Non-anonymous submissions will be rejected without review.
\author{\name Dương Mai Thương \email \\
      \addr \\
      \AND
      \name Nguyễn Phương Huy \email \\
      \addr }

\begin{document}

\renewcommand{\figurename}{Figure}
\renewcommand{\tablename}{Table}
\renewcommand{\refname}{References}
\renewcommand{\abstractname}{Abstract}

\maketitle

\begin{abstract}
Chưng cất tri thức (knowledge distillation, KD) từ một mô hình thầy (teacher) dùng nhiều điện cực có thể bù lại một phần độ chính xác bị mất khi bộ giải mã SSVEP chỉ dùng ít điện cực. Một nghiên cứu trước đã báo cáo mức cải thiện như vậy cho chưng cất quan hệ xuyên kênh (cross-channel relational KD, CCKD) trên một kiến trúc duy nhất là FB-tCNN \citep{icera2026cckd}. Chưa rõ mức cải thiện này có giữ được trên các kiến trúc khác không và đại lượng nào dự đoán được nó. Bài báo đánh giá CCKD trên năm kiến trúc công khai (FB-tCNN, DeepConvNet, ShallowConvNet, EEGNet và Compact-CNN), hai bộ dữ liệu 40 lớp (Benchmark và BETA) và hai cấu hình học trò (student) dùng ba và một điện cực. Mỗi thiết lập được đánh giá trên năm cách chia subject độc lập, mỗi cách chia tám seed. Với trọng số chưng cất cố định, mức thay đổi trung bình nằm trong khoảng $-1.30$ đến $+1.19$ điểm phần trăm (pp). Trong 20 tổ hợp kiến trúc--dữ liệu--cấu hình, tám tổ hợp cải thiện có ý nghĩa, bốn tổ hợp suy giảm có ý nghĩa và tám tổ hợp không phát hiện được hiệu ứng (kiểm định $t$ một mẫu qua các cách chia, $p<0.05$). \tbd{kết quả sau khi chọn trọng số theo validation riêng cho từng kiến trúc và thống nhất tầng chưng cất (N2a, N2b)}. Mức cải thiện đi cùng mức CCKD làm tăng độ tương đồng CKA tuyến tính giữa đặc trưng của teacher và student (Spearman $\rho=0.63$, $p=0.003$, $n=20$), trong khi độ tương đồng trước chưng cất không dự đoán được mức cải thiện ($\rho=0.06$, $p=0.80$). \tbd{kết quả thí nghiệm can thiệp theo tầng chưng cất (N5)}. Đánh giá trên một cách chia subject cho kết luận khác với năm cách chia ở 5 trên 20 tổ hợp, theo cả hai chiều. Nhiễu không xác định giữa các lần chạy đạt $0.51$~pp trên mức cải thiện, cùng bậc với nhiều hiệu ứng được báo cáo. \tbd{kết quả lượng tử hóa INT8 và fixed-point (H1, H2)}. Các kết quả cho thấy chưng cất quan hệ xuyên cấu hình điện cực phụ thuộc vào kiến trúc. Các mức cải thiện KD dưới 1--2~pp trên SSVEP cần được đánh giá trên nhiều cách chia subject.
\end{abstract}

\textbf{Keywords:} Brain--computer interface (BCI); Steady-state visual evoked potential (SSVEP); Knowledge distillation; Relational knowledge distillation; Electrode reduction; Representation similarity; Evaluation reliability.

% ======================================================================
\section{Giới thiệu}
\label{sec:introduction}

Giao diện não--máy (brain--computer interface, BCI) dựa trên điện não đồ (electroencephalography, EEG) cho phép truyền lệnh từ não tới thiết bị bên ngoài mà không cần cử động cơ. Trong các loại BCI dựa trên EEG, điện thế gợi thị giác trạng thái ổn định (steady-state visual evoked potential, SSVEP) được dùng rộng rãi nhờ tốc độ truyền thông tin cao và thời gian hiệu chuẩn ngắn \citep{chen2015highspeed,zhang2023wearable}. Các bộ nhận dạng SSVEP có độ chính xác cao thường dùng tín hiệu từ nhiều điện cực vùng đỉnh--chẩm \citep{wang2017benchmark,liu2020beta}. Số điện cực lớn làm tăng thời gian chuẩn bị và khó tích hợp vào thiết bị đeo, nên giảm số điện cực là một yêu cầu thực tế \citep{zhang2023wearable,zhu2021open,acampora2021dataset}.

Khi giảm điện cực, mô hình mất một phần thông tin không gian và độ chính xác giảm. Chưng cất tri thức (knowledge distillation, KD) là một cách bù lại: một mô hình thầy (teacher) huấn luyện trên đủ điện cực hướng dẫn một mô hình học trò (student) chỉ thấy ít điện cực \citep{hinton2015distilling,gou2021survey}. Trường hợp này khác KD thông thường ở chỗ teacher và student không quan sát cùng một đầu vào. Nếu student mô phỏng phân bố đầu ra hoặc đặc trưng tuyệt đối của teacher thì nó bị ép tái tạo thông tin mà đầu vào của nó không chứa. Chưng cất quan hệ (relational KD) tránh vấn đề này bằng cách chỉ truyền cấu trúc tương đồng giữa các mẫu \citep{tung2019spkd}.

Một nghiên cứu trước đã áp dụng ý tưởng này cho SSVEP dưới tên chưng cất quan hệ xuyên kênh (CCKD) \citep{icera2026cckd}. Trên kiến trúc FB-tCNN, CCKD cải thiện student ba điện cực và một điện cực ở cả bốn thiết lập của hai bộ dữ liệu Benchmark và BETA, với mức cải thiện trung bình từ $+0.71$ đến $+1.19$~pp qua năm cách chia subject \citep{icera2026cckd}. Nghiên cứu đó cũng nêu ba câu hỏi còn mở: mức cải thiện phụ thuộc vào tính chất nào của kiến trúc, nó thay đổi thế nào theo tập điện cực và trọng số chưng cất nên được chọn trên dữ liệu validation ra sao.

Bài báo này trả lời các câu hỏi đó. Chúng tôi không đề xuất phương pháp mới. Thay vào đó, chúng tôi đánh giá có kiểm soát một phương pháp đã biết trên năm kiến trúc và tìm điều kiện để nó có ích. Bài báo xét bốn câu hỏi nghiên cứu (research question, RQ):

\begin{itemize}
    \item \textbf{RQ1.} Mức cải thiện của CCKD có giữ được khi thay kiến trúc và khi thay tập điện cực của student không?
    \item \textbf{RQ2.} Đại lượng nào dự đoán được mức cải thiện? Cụ thể, độ tương đồng biểu diễn giữa teacher và student có đóng vai trò này không?
    \item \textbf{RQ3.} Đánh giá trên một cách chia subject có cho cùng kết luận với đánh giá trên nhiều cách chia không?
    \item \textbf{RQ4.} Mức cải thiện có còn khi mô hình chạy ở độ chính xác số học thấp (INT8 trên CPU, fixed-point trên FPGA) không?
\end{itemize}

Các đóng góp chính của bài báo gồm:

\begin{itemize}
    \item Một đánh giá có kiểm soát của CCKD trên năm kiến trúc, hai bộ dữ liệu, hai số điện cực và \tbd{số} tập điện cực, mỗi thiết lập dùng năm cách chia subject và tám seed. Tầng chưng cất được định nghĩa thống nhất cho mọi kiến trúc và trọng số chưng cất được chọn trên validation riêng cho từng kiến trúc (Mục~\ref{sec:rq1}).
    \item Một phân tích cho thấy mức tăng độ tương đồng CKA do CCKD tạo ra đi cùng mức cải thiện độ chính xác, kèm thí nghiệm can thiệp theo tầng chưng cất để kiểm tra quan hệ này (Mục~\ref{sec:cka}).
    \item Một phân tích độ tin cậy của đánh giá: so sánh kết luận từ một và năm cách chia subject, đo nhiễu không xác định giữa các lần chạy và kiểm tra việc chọn siêu tham số trên một cách chia (Mục~\ref{sec:rq3}).
    \item Kết quả về mức cải thiện sau lượng tử hóa INT8 và fixed-point, kèm độ trễ đo trên Raspberry Pi 4 và ZCU104 (Mục~\ref{sec:rq4}).
\end{itemize}

Phần còn lại của bài được tổ chức như sau. Mục~\ref{sec:related} tổng quan các nghiên cứu liên quan. Mục~\ref{sec:method} tóm tắt CCKD. Mục~\ref{sec:setup} trình bày thiết lập thực nghiệm. Các Mục~\ref{sec:rq1} đến~\ref{sec:rq4} lần lượt trả lời RQ1 đến RQ4. Mục~\ref{sec:discussion} thảo luận và nêu giới hạn. Mục~\ref{sec:conclusion} kết luận.

% ======================================================================
\section{Nghiên cứu liên quan}
\label{sec:related}

\paragraph{Học sâu cho nhận dạng SSVEP.} Các phương pháp không cần huấn luyện như phân tích tương quan chính tắc (canonical correlation analysis, CCA) và FBCCA có chi phí thấp nhưng giảm hiệu năng khi cửa sổ tín hiệu ngắn \citep{lin2007cca,chen2015fbcca}. Các mô hình học sâu học trực tiếp đặc trưng không gian và thời gian từ EEG. Các kiến trúc tiêu biểu gồm EEGNet, Compact-CNN, FB-tCNN, DeepConvNet, ShallowConvNet và các mô hình dựa trên Transformer \citep{lawhern2018eegnet,waytowich2018compact,ding2021filterbank,schirrmeister2017deep,guney2023deep,chen2023transformer,wu2024analysis}. Các kiến trúc này khác nhau ở cách gộp chiều điện cực: tích chập không gian thường (ordinary spatial convolution) ở FB-tCNN, DeepConvNet và ShallowConvNet, tích chập theo chiều sâu (depthwise convolution) ở EEGNet và Compact-CNN. Bài báo dùng chính sự khác biệt này làm một biến giải thích ở RQ2.

\paragraph{Chưng cất tri thức khi teacher và student có đầu vào khác nhau.} KD ban đầu truyền phân bố đầu ra đã làm mềm của teacher \citep{hinton2015distilling}. Các biến thể sau truyền đặc trưng trung gian (FitNets), bản đồ chú ý (attention transfer) hoặc quan hệ giữa các mẫu \citep{romero2015fitnets,zagoruyko2017paying,tung2019spkd,gou2021survey,moslemi2024survey}. Trong EEG, KD đã được dùng cho nhận dạng cảm xúc và nén mô hình \citep{wang2024kdemotion}. Một số nghiên cứu SSVEP xét trường hợp teacher và student thấy đầu vào khác nhau: RMKD chưng cất từ cửa sổ dài sang cửa sổ ngắn \citep{lan2025rmkd}, MESD chưng cất từ EEG da đầu sang EEG tai bằng cách kết hợp chưng cất đặc trưng và chưng cất đầu ra \citep{sun2024mesd}. CCKD xét trường hợp giảm điện cực trên da đầu và chỉ dùng chưng cất quan hệ \citep{icera2026cckd}. Các nghiên cứu này đều báo cáo trên một kiến trúc chính, nên chưa cho biết hiệu quả của KD xuyên đầu vào phụ thuộc vào kiến trúc đến mức nào.

\paragraph{Điều kiện để KD có hiệu quả.} Nhiều nghiên cứu trong học máy chỉ ra rằng student không phải lúc nào cũng khớp được teacher và mức cải thiện của KD phụ thuộc vào khoảng cách năng lực giữa hai mô hình \cantrich{nghiên cứu về khoảng cách năng lực teacher--student và độ trung thành của KD, Q1/Q2 từ 2020}. Bài báo này xét một nguồn chênh lệch khác: chênh lệch thông tin đầu vào do teacher và student thấy các tập điện cực khác nhau.

\paragraph{Đo độ tương đồng biểu diễn.} CKA (centered kernel alignment) đo độ tương đồng giữa hai tập đặc trưng của cùng các mẫu, bất biến với phép quay trực giao và phép co giãn đẳng hướng \citep{kornblith2019similarity}. Hàm mất mát quan hệ của CCKD so sánh ma trận Gram của teacher và student. CKA tuyến tính cũng được tính từ ma trận Gram, nên nó là thước đo tự nhiên để kiểm tra CCKD đã thay đổi biểu diễn của student đến đâu.

\paragraph{Độ tin cậy của đánh giá cross-subject.} Nhiều nghiên cứu SSVEP báo cáo kết quả trên một cách chia subject duy nhất hoặc dùng leave-one-subject-out \cantrich{khảo sát về giao thức đánh giá cross-subject trong EEG/SSVEP, Q1/Q2 từ 2020}. Khi mức cải thiện chỉ khoảng 1~pp và tập kiểm thử chỉ có 10--20 subject, kết luận có thể phụ thuộc vào cách chia. Bài báo đo trực tiếp mức phụ thuộc này ở RQ3.

% ======================================================================
\section{Chưng cất quan hệ xuyên kênh (CCKD)}
\label{sec:method}

Mục này tóm tắt CCKD theo \citet{icera2026cckd} và nêu các thay đổi về giao thức trong bài báo này. Teacher nhận tín hiệu từ chín điện cực đỉnh--chẩm và được huấn luyện trước bằng cross-entropy rồi giữ cố định. Student có cùng kiến trúc với teacher nhưng chỉ nhận ba hoặc một điện cực.

Với một mini-batch $N$ mẫu, gọi $F_T^{(l)} \in \mathbb{R}^{N \times d_T}$ và $F_S^{(l)} \in \mathbb{R}^{N \times d_S}$ là đặc trưng của teacher và student tại tầng $l$. Ma trận Gram của mỗi mô hình là
\begin{equation}
G_{T}^{(l)} = F_{T}^{(l)} \bigl(F_{T}^{(l)}\bigr)^{\top}, \qquad
G_{S}^{(l)} = F_{S}^{(l)} \bigl(F_{S}^{(l)}\bigr)^{\top} \in \mathbb{R}^{N \times N}.
\label{eq:gram}
\end{equation}
Ma trận Gram luôn có kích thước $N \times N$, nên CCKD không đòi hỏi $d_T = d_S$. Mỗi hàng được chuẩn hóa theo chuẩn $\ell_2$:
\begin{equation}
\hat{G}^{(l)}[i,:] = \frac{G^{(l)}[i,:]}{\lVert G^{(l)}[i,:] \rVert_2}, \qquad i = 1, \dots, N.
\label{eq:norm}
\end{equation}
Hàm mất mát quan hệ tại tầng $l$ là
\begin{equation}
\mathcal{L}_{\text{CCKD}}^{(l)} = \frac{1}{N^2} \bigl\lVert \hat{G}_{T}^{(l)} - \hat{G}_{S}^{(l)} \bigr\rVert_F^2 .
\label{eq:relloss}
\end{equation}

\subsection{Hàm mục tiêu và tầng chưng cất}
\label{sec:training}

Với tập tầng chưng cất $\mathcal{K}$, student tối ưu
\begin{equation}
\mathcal{L} = \mathcal{L}_{\text{CE}} + \lambda_{sp} \, \frac{1}{\lvert \mathcal{K} \rvert} \sum_{l \in \mathcal{K}} \mathcal{L}_{\text{CCKD}}^{(l)},
\label{eq:total_loss}
\end{equation}
trong đó $\mathcal{L}_{\text{CE}}$ là cross-entropy của student và $\lambda_{sp}$ là trọng số chưng cất. Sau huấn luyện, teacher bị loại bỏ, nên CCKD không làm thay đổi kiến trúc, số tham số hay chi phí suy luận của student.

Trong nghiên cứu trước, FB-tCNN được chưng cất ở hai tầng, còn khi mở rộng sang các kiến trúc khác thì chỉ tầng cuối được chưng cất. Như vậy kiến trúc và số tầng chưng cất bị lẫn vào nhau. Bài báo này định nghĩa hai tầng giống nhau cho mọi kiến trúc:
\begin{itemize}
    \item \textbf{Tầng giữa} (\texttt{mid}): đầu ra của khối tích chập không gian ngay sau chuẩn hóa batch và hàm phi tuyến, trước pooling và dropout. Đây là điểm chiều điện cực vừa bị gộp, nên độ rộng của tầng này không phụ thuộc vào số điện cực đầu vào.
    \item \textbf{Tầng cuối} (\texttt{last}): đặc trưng áp chót, tức đầu vào của bộ phân loại.
\end{itemize}
Giao thức chính chưng cất chỉ tầng cuối ($\mathcal{K}=\{\texttt{last}\}$). Thí nghiệm can thiệp ở RQ2 so sánh thêm $\mathcal{K}=\{\texttt{mid}\}$ và $\mathcal{K}=\{\texttt{mid},\texttt{last}\}$. Kết quả theo giao thức cũ được ghi rõ là ``giao thức cũ'' mỗi khi được dùng.

% ======================================================================
\section{Thiết lập thực nghiệm}
\label{sec:setup}

\subsection{Dữ liệu và tập điện cực}

Bài báo dùng hai bộ dữ liệu SSVEP công khai 40 lớp, lấy mẫu 250~Hz: Benchmark gồm 35 subject \citep{wang2017benchmark} và BETA gồm 70 subject \citep{liu2020beta}. Thiết lập được tóm tắt trong Bảng~\ref{tab:settings}. Mỗi cách chia subject gán các subject rời nhau vào tập huấn luyện, validation và kiểm thử. Tập validation dùng để chọn checkpoint và siêu tham số. Tập kiểm thử chỉ được đánh giá một lần. Năm cách chia dùng các seed chia 2026, 3026, 4026, 5026 và 6026.

\begin{table}[h]
\caption{Summary of experimental settings.}
\label{tab:settings}
\begin{center}
\begin{tabular}{lcccccc}
\toprule
Dataset & Subjects & Train/Val/Test & Splits & Window & Teacher & Student montages \\
\midrule
Benchmark & 35 & 21 / 4 / 10 & 5 & 1.0 s & 9 channels & 3C, 1C \tbd{+ alternatives} \\
BETA      & 70 & 42 / 8 / 20 & 5 & 1.0 s & 9 channels & 3C, 1C \tbd{+ alternatives} \\
\bottomrule
\end{tabular}
\end{center}
\end{table}

Teacher dùng chín điện cực Pz, PO5, PO3, POz, PO4, PO6, O1, Oz và O2. Tập điện cực gốc của student là O1, Oz, O2 (3C) và Oz (1C), giống nghiên cứu trước. Để trả lời RQ1 về tập điện cực, bài báo thêm hai tập ba điện cực là \{PO3, POz, PO4\} và \{POz, O1, O2\}, cùng hai tập một điện cực là \{POz\} và \{O1\}. Các tập này được cố định trước khi chạy thí nghiệm.

\subsection{Năm kiến trúc}
\label{sec:backbones}

Bảng~\ref{tab:backbones} liệt kê năm kiến trúc cùng các tính chất được dùng làm biến giải thích ở RQ2. Một số siêu tham số kiến trúc được điều chỉnh cho cửa sổ 1.0~s (250 mẫu). DeepConvNet giữ ba khối tích chập--pooling thay vì bốn. ShallowConvNet dùng pooling 35 mẫu, bước 7. Compact-CNN dùng kernel thời gian 64 thay vì 256.

\begin{table}[h]
\caption{Backbones and their architectural properties. ``mid'' and ``last'' are the widths of the two distillation levels; ``mid'' does not depend on the number of input electrodes.}
\label{tab:backbones}
\begin{center}
\resizebox{\linewidth}{!}{%
\begin{tabular}{lcccccc}
\toprule
Backbone & Spatial conv & Nonlinearity & Max temporal kernel & mid & last & Params (9C) \\
\midrule
FB-tCNN \citep{ding2021filterbank} & ordinary, shared across 3 bands & ELU & 250 & 12000 & 32 & 90,536 \\
DeepConvNet \citep{schirrmeister2017deep} & ordinary & ELU & 10 & 6250 & 900 & \tbd{} \\
ShallowConvNet \citep{schirrmeister2017deep} & ordinary & square, log & 25 & 10000 & 1240 & \tbd{} \\
EEGNet \citep{lawhern2018eegnet} & depthwise & ELU & 64 & 4016 & 112 & \tbd{} \\
Compact-CNN \citep{waytowich2018compact} & depthwise & ELU & 64 & 24096 & 672 & 45,256 \\
\bottomrule
\end{tabular}}
\end{center}
\end{table}

\subsection{Huấn luyện và chọn siêu tham số}

Mọi mô hình được huấn luyện 100 epoch bằng Adam (tốc độ học $10^{-3}$, weight decay $10^{-4}$), lịch cosine annealing và batch 64. Checkpoint được chọn theo độ chính xác validation trung bình theo subject. Mỗi cấu hình chạy tám seed khởi tạo (2026--2033) trên mỗi cách chia. Các seed của Python, NumPy và PyTorch được cố định.

Trọng số chưng cất được chọn trong $\lambda_{sp} \in \{1, 10, 100, 1000\}$ theo độ chính xác validation, riêng cho từng kiến trúc, bộ dữ liệu, cấu hình và cách chia. Tập kiểm thử không tham gia vào lựa chọn này. Kết quả với $\lambda_{sp}=100$ cố định (giá trị của nghiên cứu trước) được báo cáo song song để so sánh.

\subsection{Thước đo và kiểm định}

Thước đo chính là độ chính xác phân loại trên các subject kiểm thử. Mức cải thiện $\Delta$ là chênh lệch theo cặp giữa student chưng cất và student baseline (chỉ huấn luyện bằng cross-entropy) trong cùng lần chạy, tính theo điểm phần trăm (pp), lấy trung bình qua subject sau khi lấy trung bình qua seed. Trên một cách chia, ý nghĩa thống kê được kiểm bằng Wilcoxon signed-rank qua subject \citep{wilcoxon1945individual}. Qua năm cách chia, ý nghĩa được kiểm bằng kiểm định $t$ một mẫu trên năm giá trị $\Delta$ ($H_0$: $\bar\Delta = 0$). Khi nhiều kiểm định thuộc cùng một họ, giá trị $p$ được hiệu chỉnh Holm--Bonferroni \cantrich{Holm (1979)}. Kết luận ``không có hiệu ứng'' chỉ được đưa ra khi kiểm định tương đương TOST đạt ý nghĩa \cantrich{Lakens, TOST}.

\subsection{Hạn chế kỹ thuật đã biết}

Ba hạn chế kỹ thuật ảnh hưởng tới cách đọc kết quả. Nếu không nêu rõ thì người đọc có thể hiểu sai độ chính xác của các con số.
\begin{itemize}
    \item Huấn luyện GPU không tái lập đúng từng bit. Mức nhiễu này được đo ở Mục~\ref{sec:nondeterminism}.
    \item Student baseline dùng độ chính xác hỗn hợp (automatic mixed precision, AMP), còn student chưng cất dùng FP32. Thiết kế theo cặp giữ cùng chênh lệch này ở mọi kiến trúc, nhưng chúng tôi chưa đo riêng ảnh hưởng của nó.
    \item Trong một seed, teacher, baseline và các student chưng cất được huấn luyện nối tiếp từ cùng một trạng thái sinh số ngẫu nhiên. Vì vậy cùng nhãn seed không bảo đảm cùng khởi tạo giữa các phương pháp.
\end{itemize}

% ======================================================================
\section{RQ1: Mức cải thiện theo kiến trúc và tập điện cực}
\label{sec:rq1}
\label{sec:results}
\label{sec:generalization}

\subsection{Kết quả theo giao thức cũ}

Bảng~\ref{tab:five_backbones} tổng hợp mức cải thiện của CCKD trên năm kiến trúc theo giao thức cũ. Theo giao thức này, $\lambda_{sp}=100$ được giữ cố định cho mọi kiến trúc, FB-tCNN được chưng cất ở hai tầng và các kiến trúc khác chỉ ở tầng cuối. Dòng FB-tCNN lấy từ nghiên cứu trước \citep{icera2026cckd}.

\begin{table}[h]
\caption{CCKD gain over five subject splits under the earlier protocol ($\lambda_{sp}=100$ for all backbones; FB-tCNN distills two levels, the other backbones the last level only). Entries: mean $\Delta$ (pp) and one-sample $t$-test $p$ over splits. Bold: $p<0.05$. FB-tCNN row from \citet{icera2026cckd}.}
\label{tab:five_backbones}
\begin{center}
\begin{tabular}{lcccc}
\toprule
Backbone & Benchmark 3C & Benchmark 1C & BETA 3C & BETA 1C \\
\midrule
FB-tCNN & \textbf{+1.00 (0.0050)} & \textbf{+0.71 (0.0103)} & \textbf{+1.19 (0.0002)} & \textbf{+0.72 (0.0293)} \\
DeepConvNet & \textbf{+0.59 (0.0006)} & +0.13 (0.684) & \textbf{+0.84 (0.0047)} & \textbf{+0.61 (0.0161)} \\
ShallowConvNet & \textbf{$-0.91$ (0.0125)} & \textbf{$-0.95$ (0.0161)} & +0.46 (0.191) & +0.17 (0.300) \\
EEGNet & \textbf{+1.09 (0.0179)} & $-0.66$ (0.120) & \textbf{$-1.30$ (0.0006)} & \textbf{$-1.29$ (0.0036)} \\
Compact-CNN & +0.06 (0.872) & $-0.27$ (0.128) & $-0.24$ (0.190) & $-0.20$ (0.219) \\
\bottomrule
\end{tabular}
\end{center}
\end{table}

Trong 20 tổ hợp, tám tổ hợp cải thiện có ý nghĩa, bốn tổ hợp suy giảm có ý nghĩa và tám tổ hợp không đạt ý nghĩa (Bảng~\ref{tab:five_backbones}). Kết quả không xếp được theo nhóm kiến trúc. Trong nhóm tích chập không gian thường, FB-tCNN cải thiện ở cả bốn thiết lập, DeepConvNet cải thiện ở ba thiết lập, còn ShallowConvNet suy giảm ở cả hai thiết lập của Benchmark. Trong nhóm tích chập theo chiều sâu, EEGNet cải thiện ở Benchmark-3C nhưng suy giảm ở cả hai thiết lập của BETA, còn Compact-CNN không cho hiệu ứng nào đạt ý nghĩa. Các tổ hợp không đạt ý nghĩa được xếp là ``không phát hiện được hiệu ứng'', chưa phải ``không có hiệu ứng''. \tbd{TOST trên dữ liệu năm cách chia; phân tích trước chỉ chạy TOST trên một cách chia}.

Bảng~\ref{tab:five_backbones} chưa trả lời được RQ1, vì hai yếu tố bị lẫn vào biến kiến trúc. Thứ nhất, $\lambda_{sp}=100$ được chọn cho FB-tCNN. Nếu không chọn lại thì một kiến trúc có thể suy giảm chỉ vì trọng số không phù hợp với nó. Thứ hai, số tầng chưng cất khác nhau giữa FB-tCNN và các kiến trúc còn lại. Mục~\ref{sec:rq1_matched} loại bỏ cả hai yếu tố này.

\subsection{Kết quả với tầng chưng cất thống nhất và trọng số chọn theo validation}
\label{sec:rq1_matched}

\tbd{Bảng N2: 5 kiến trúc $\times$ 4 thiết lập, chưng cất tầng cuối; hai khối dòng: $\lambda_{sp}=100$ cố định và $\lambda_{sp}$ chọn theo validation; mỗi ô $\bar\Delta$ ($p_t$, số cách chia dương); cột $p_{\text{Holm}}$ trên họ 20 kiểm định. Nguồn: \texttt{results/module3/n2a}, \texttt{n2b}.}

\tbd{Đoạn diễn giải: số tổ hợp đổi kết luận so với Bảng~\ref{tab:five_backbones}; phân bố $\lambda_{sp}$ được chọn theo kiến trúc.}

\subsection{Độ nhạy theo tập điện cực}
\label{sec:montage}

\tbd{Bảng N3: 5 kiến trúc $\times$ 2 bộ dữ liệu $\times$ 6 tập điện cực (3C, 3C\_PPO, 3C\_POO, 1C, 1C\_POz, 1C\_O1), $\lambda_{sp}$ theo N2b; nguồn \texttt{results/module3/n3}.}

% ======================================================================
\section{RQ2: Đại lượng dự đoán mức cải thiện}
\label{sec:rq2}

\subsection{Độ tương đồng CKA giữa teacher và student}
\label{sec:cka}

Hàm mất mát~\eqref{eq:relloss} kéo ma trận Gram của student về phía ma trận Gram của teacher. Nếu CCKD có ích bằng cơ chế này thì mức cải thiện độ chính xác phải đi cùng mức tăng tương đồng biểu diễn. Chúng tôi đo độ tương đồng bằng CKA tuyến tính \citep{kornblith2019similarity} giữa đặc trưng tầng cuối của teacher và student trên tập kiểm thử. Ký hiệu $\Delta$CKA là CKA của student chưng cất trừ CKA của student baseline.

Phân tích sơ bộ dùng một seed trên cách chia 2026 theo giao thức cũ (Bảng~\ref{tab:cka}). $\Delta$CKA dương ở hầu hết các tổ hợp. Ngoại lệ là ShallowConvNet, có $\Delta$CKA xấp xỉ 0 ở cả hai bộ dữ liệu, cũng là kiến trúc duy nhất suy giảm ổn định trên Benchmark. Trên 20 tổ hợp, $\Delta$CKA tương quan với $\bar\Delta$ qua năm cách chia (Pearson $r=0.58$, $p=0.007$; Spearman $\rho=0.63$, $p=0.003$). Khi bỏ lần lượt từng kiến trúc, giá trị $p$ của tương quan vẫn dưới $0.06$ ở mọi trường hợp. CKA của student baseline không tương quan với mức cải thiện ($r=0.13$, $p=0.58$; $\rho=0.06$, $p=0.80$). Như vậy, mức chênh biểu diễn ban đầu giữa teacher và student không dự đoán được hiệu quả của CCKD. Điều dự đoán được là mức CCKD thu hẹp chênh lệch đó.

\begin{table}[h]
\caption{Linear CKA between teacher and student last-level features before and after CCKD (preliminary: one seed, split 2026, earlier protocol; 3C and 1C averaged per backbone and dataset).}
\label{tab:cka}
\begin{center}
\begin{tabular}{llccc}
\toprule
Backbone & Dataset & CKA (baseline) & CKA (CCKD) & $\Delta$CKA \\
\midrule
\multirow{2}{*}{FB-tCNN} & Benchmark & 0.45 & 0.51 & +0.06 \\
 & BETA & 0.14 & 0.17 & +0.03 \\
\multirow{2}{*}{DeepConvNet} & Benchmark & 0.51 & 0.55 & +0.04 \\
 & BETA & 0.23 & 0.26 & +0.03 \\
\multirow{2}{*}{ShallowConvNet} & Benchmark & 0.67 & 0.67 & $-0.01$ \\
 & BETA & 0.38 & 0.38 & +0.00 \\
\multirow{2}{*}{EEGNet} & Benchmark & 0.43 & 0.50 & +0.08 \\
 & BETA & 0.06 & 0.07 & +0.01 \\
\multirow{2}{*}{Compact-CNN} & Benchmark & 0.49 & 0.50 & +0.01 \\
 & BETA & 0.02 & 0.03 & +0.00 \\
\bottomrule
\end{tabular}
\end{center}
\end{table}

Phân tích này có ba giới hạn. Nó chỉ dùng một seed, lẫn yếu tố số tầng chưng cất của giao thức cũ và chỉ cho thấy tương quan. Trường hợp EEGNet--BETA có $\Delta$CKA gần 0 nhưng suy giảm mạnh nhất ($-1.30$~pp), nên còn yếu tố khác chưa được nắm bắt. Các thí nghiệm mới tính CKA ở cả hai tầng cho mọi lần chạy, trên năm cách chia và tám seed.

\tbd{Bảng N4: tương quan $\Delta$CKA--$\bar\Delta$ ở tầng giữa và tầng cuối, 8 seed $\times$ 5 cách chia, giao thức thống nhất; khoảng tin cậy bootstrap.}

\subsection{Thí nghiệm can thiệp theo tầng chưng cất}
\label{sec:layer_intervention}

Tương quan ở trên không phân biệt được hai giải thích. Hoặc mức tăng tương đồng gây ra cải thiện, hoặc cả hai cùng do một tính chất của kiến trúc. Để tách hai khả năng, chúng tôi thay đổi tầng chưng cất trong khi giữ nguyên kiến trúc, dữ liệu và cách chia.

\tbd{Bảng N5: $\bar\Delta$ và $\Delta$CKA theo $\mathcal{K} \in \{\texttt{last}, \texttt{mid}, \texttt{mid+last}\}$ cho 5 kiến trúc $\times$ 4 thiết lập.}

% ======================================================================
\section{RQ3: Độ tin cậy của đánh giá}
\label{sec:rq3}
\label{sec:mc_multisplit}

\subsection{Một cách chia so với năm cách chia}

Bảng~\ref{tab:single_vs_multi} so sánh kết luận rút ra từ cách chia 2026 với kết luận từ năm cách chia cho 20 tổ hợp của Bảng~\ref{tab:five_backbones}. Kết luận trên một cách chia dựa vào kiểm định Wilcoxon qua subject, còn kết luận trên năm cách chia dựa vào kiểm định $t$ qua cách chia. Mỗi kết luận được xếp vào một trong ba loại: cải thiện ($+$), suy giảm ($-$) hoặc không đạt ý nghĩa ($0$), với ngưỡng $p<0.05$.

\begin{table}[h]
\caption{Verdict from split 2026 alone (Wilcoxon over test subjects, 8 seeds) vs.\ five splits (one-sample $t$-test over split-level gains) for CCKD under the earlier protocol. $+$/$-$: significant gain/loss ($p<0.05$); 0: not significant. Rows where the verdicts differ are marked with $\star$. Source: \texttt{scripts/single\_vs\_multisplit.py}.}
\label{tab:single_vs_multi}
\begin{center}
\footnotesize
\begin{tabular}{lllcccccl}
\toprule
Backbone & Dataset & Student & $\Delta_{2026}$ (pp) & $p_{\text{W}}$ & $\bar\Delta_{5}$ (pp) & Pos.\ splits & $p_t$ & Verdict (1 $\to$ 5) \\
\midrule
FB-tCNN & Benchmark & 3C & +0.45 & 0.375 & +1.00 & 5/5 & 0.0050 & 0 $\to$ $+$ $\star$ \\
 & & 1C & +0.59 & 0.193 & +0.71 & 5/5 & 0.0103 & 0 $\to$ $+$ $\star$ \\
 & BETA & 3C & +1.05 & 0.012 & +1.19 & 5/5 & 0.0002 & $+$ $\to$ $+$ \\
 & & 1C & +1.06 & 0.004 & +0.72 & 4/5 & 0.0293 & $+$ $\to$ $+$ \\
DeepConvNet & Benchmark & 3C & +0.52 & 0.012 & +0.59 & 5/5 & 0.0006 & $+$ $\to$ $+$ \\
 & & 1C & +1.17 & 0.084 & +0.13 & 2/5 & 0.684 & 0 $\to$ 0 \\
 & BETA & 3C & +0.90 & 0.006 & +0.84 & 5/5 & 0.0047 & $+$ $\to$ $+$ \\
 & & 1C & +0.56 & 0.023 & +0.61 & 5/5 & 0.0161 & $+$ $\to$ $+$ \\
ShallowConvNet & Benchmark & 3C & $-0.31$ & 0.432 & $-0.91$ & 0/5 & 0.0125 & 0 $\to$ $-$ $\star$ \\
 & & 1C & $-0.55$ & 0.016 & $-0.95$ & 0/5 & 0.0161 & $-$ $\to$ $-$ \\
 & BETA & 3C & +0.50 & 0.040 & +0.46 & 4/5 & 0.191 & $+$ $\to$ 0 $\star$ \\
 & & 1C & +0.32 & 0.239 & +0.17 & 3/5 & 0.300 & 0 $\to$ 0 \\
EEGNet & Benchmark & 3C & +0.66 & 0.084 & +1.09 & 5/5 & 0.0179 & 0 $\to$ $+$ $\star$ \\
 & & 1C & $-0.17$ & 0.492 & $-0.66$ & 0/5 & 0.120 & 0 $\to$ 0 \\
 & BETA & 3C & $-1.30$ & 0.003 & $-1.29$ & 0/5 & 0.0006 & $-$ $\to$ $-$ \\
 & & 1C & $-1.02$ & 0.004 & $-1.29$ & 0/5 & 0.0036 & $-$ $\to$ $-$ \\
Compact-CNN & Benchmark & 3C & $-0.21$ & 0.742 & +0.06 & 1/5 & 0.872 & 0 $\to$ 0 \\
 & & 1C & $-0.52$ & 0.375 & $-0.27$ & 1/5 & 0.128 & 0 $\to$ 0 \\
 & BETA & 3C & $-0.04$ & 0.695 & $-0.24$ & 1/5 & 0.190 & 0 $\to$ 0 \\
 & & 1C & $-0.12$ & 0.526 & $-0.20$ & 1/5 & 0.219 & 0 $\to$ 0 \\
\bottomrule
\end{tabular}
\end{center}
\end{table}

Kết luận thay đổi ở 5 trên 20 tổ hợp, theo cả hai chiều (Bảng~\ref{tab:single_vs_multi}). Ở ba tổ hợp, một cách chia bỏ sót một cải thiện ổn định: FB-tCNN Benchmark-3C và 1C, EEGNet Benchmark-3C. Ở một tổ hợp, một cách chia bỏ sót một suy giảm ổn định: ShallowConvNet Benchmark-3C. Ở một tổ hợp, một cách chia báo cáo một cải thiện không lặp lại: ShallowConvNet BETA-3C. Ngoài ra, DeepConvNet Benchmark-1C cho $+1.17$~pp trên cách chia 2026 nhưng chỉ $+0.13$~pp trung bình qua năm cách chia. Các sai lệch tập trung ở Benchmark, nơi tập kiểm thử chỉ có 10 subject.

Hai kiểm định không cùng đơn vị lặp: Wilcoxon lặp qua subject, kiểm định $t$ lặp qua cách chia. Vì vậy bảng không cho biết kiểm định nào có công suất cao hơn. Điều bảng cho thấy là một nghiên cứu chỉ báo cáo một cách chia có thể đi đến kết luận ngược với kết luận khi lặp lại trên các tập subject khác.

\subsection{Nhiễu không xác định giữa các lần chạy}
\label{sec:nondeterminism}

Chúng tôi lặp lại năm lần một cấu hình giống hệt nhau (FB-tCNN, $\lambda_{sp}=100$, seed 2026, cách chia 2026) theo hai cách. Cách thứ nhất chạy năm tiến trình riêng. Cách thứ hai lặp năm lần trong cùng một tiến trình. Bảng~\ref{tab:nondeterminism} cho thấy teacher tái lập gần như đúng từng bit (SD $\le 0.02$~pp). Với các student, độ lệch chuẩn của $\Delta$ dao động từ $0.04$ đến $0.51$~pp, lớn nhất ở cấu hình một điện cực. Nhiễu cỡ $0.3$--$0.5$~pp cùng bậc với nhiều hiệu ứng trong bài, nên các so sánh có $\lvert\Delta\rvert < 0.5$~pp chỉ có nghĩa khi đã lấy trung bình qua nhiều seed và nhiều cách chia.

\begin{table}[h]
\caption{Run-to-run nondeterminism over five repetitions of an identical configuration (FB-tCNN, $\lambda_{sp}=100$, seed 2026, split 2026). ``Across processes'': five separate script calls. ``Within process'': five repetitions in one Python/CUDA process. Units: pp.}
\label{tab:nondeterminism}
\begin{center}
\footnotesize
\begin{tabular}{llcccc}
\toprule
Dataset & Quantity & \multicolumn{2}{c}{Across processes} & \multicolumn{2}{c}{Within process} \\
 & & range & SD & range & SD \\
\midrule
\multirow{7}{*}{Benchmark} & Teacher-9C & 0.00 & 0.00 & 0.00 & 0.00 \\
 & 3C baseline & 0.42 & 0.19 & 0.04 & 0.02 \\
 & 1C baseline & 0.46 & 0.18 & 0.25 & 0.11 \\
 & 3C CCKD & 0.17 & 0.06 & 0.13 & 0.05 \\
 & 1C CCKD & 0.79 & 0.40 & 1.04 & 0.40 \\
 & $\Delta$ (3C) & 0.46 & 0.19 & 0.08 & 0.04 \\
 & $\Delta$ (1C) & 1.08 & 0.51 & 0.92 & 0.33 \\
\midrule
\multirow{7}{*}{BETA} & Teacher-9C & 0.03 & 0.01 & 0.03 & 0.02 \\
 & 3C baseline & 0.09 & 0.05 & 0.09 & 0.04 \\
 & 1C baseline & 0.06 & 0.02 & 0.31 & 0.12 \\
 & 3C CCKD & 0.19 & 0.08 & 0.25 & 0.10 \\
 & 1C CCKD & 0.13 & 0.05 & 0.50 & 0.26 \\
 & $\Delta$ (3C) & 0.19 & 0.07 & 0.22 & 0.08 \\
 & $\Delta$ (1C) & 0.16 & 0.06 & 0.78 & 0.37 \\
\bottomrule
\end{tabular}
\end{center}
\end{table}

\subsection{Chọn siêu tham số trên một cách chia}
\label{sec:lambda_val}

Chọn $\lambda_{sp}$ theo validation trên cách chia 2026 cho FB-tCNN dẫn đến $\lambda_{sp}=1$ trên Benchmark và $\lambda_{sp}=100$ trên BETA (Bảng~\ref{tab:lambda_val} ở Phụ lục). Chênh lệch validation giữa các giá trị trên Benchmark chỉ khoảng $0.2$--$0.3$~pp. Khi lặp lại trên bốn cách chia còn lại, $\lambda_{sp}=1$ không cho cải thiện ổn định, còn $\lambda_{sp}=100$ cải thiện ở cả năm cách chia (Bảng~\ref{tab:lambda_robust}). Như vậy, lựa chọn theo validation trên một cách chia có thể không lặp lại được. Đây là lý do giao thức mới chọn $\lambda_{sp}$ riêng trên từng cách chia và báo cáo phân bố của các giá trị được chọn.

\begin{table}[h]
\caption{Robustness of $\lambda_{sp}=1$ (selected on validation of split 2026) and $\lambda_{sp}=100$ over five Benchmark splits (FB-tCNN, earlier protocol, 8 seeds per split). $\Delta$ (pp) per split; $p_t$: one-sample $t$-test over the five gains.}
\label{tab:lambda_robust}
\begin{center}
\footnotesize
\begin{tabular}{llccccccc}
\toprule
$\lambda_{sp}$ & Student & 2026 & 3026 & 4026 & 5026 & 6026 & $\bar\Delta \pm$ SD & $p_t$ \\
\midrule
\multirow{2}{*}{1} & 3C & +0.22 & $-0.36$ & $-0.23$ & +0.60 & +0.52 & +0.15 $\pm$ 0.43 & 0.4845 \\
 & 1C & $-0.33$ & +0.29 & +0.01 & $-0.22$ & +0.29 & +0.01 $\pm$ 0.28 & 0.9631 \\
\multirow{2}{*}{100} & 3C & +0.44 & +1.17 & +0.94 & +0.95 & +1.49 & +1.00 $\pm$ 0.38 & \textbf{0.0043} \\
 & 1C & +0.73 & +0.59 & +0.83 & +0.74 & +0.60 & +0.70 $\pm$ 0.10 & \textbf{0.0001} \\
\bottomrule
\end{tabular}
\end{center}
\end{table}

% ======================================================================
\section{RQ4: Mức cải thiện ở độ chính xác số học thấp}
\label{sec:rq4}
\label{sec:pi4}

Nghiên cứu trước đã đo độ trễ FP32 của FB-tCNN trên Raspberry Pi 4 và thấy giảm điện cực gần như không giảm tính toán, vì tầng đầu gộp chiều điện cực \citep{icera2026cckd}. Vì vậy mục này không xét tốc độ. Câu hỏi ở đây là mức cải thiện cỡ 1~pp của CCKD có bị lượng tử hóa xóa mất hay không.

\tbd{H1: $\Delta$ trước và sau INT8 PTQ cho 5 kiến trúc $\times$ 4 thiết lập trên 5 cách chia.}

\tbd{H2: $\Delta$ theo độ rộng bit fixed-point (16, 12, 10, 8 bit) bằng mô phỏng bit-accurate hls4ml.}

\tbd{H3: độ trễ và bộ nhớ FP32/INT8 của 5 kiến trúc trên Raspberry Pi 4 (Cortex-A72) và ZCU104 (Cortex-A53), batch 1, cửa sổ 1.0 s; đường Pareto độ chính xác--độ trễ.}

\tbd{H4 (tùy chọn): một kiến trúc tổng hợp trên PL của ZCU104 (Vitis HLS 2024.2): độ trễ, LUT/DSP/BRAM, công suất, xác minh độ chính xác trên board.}

% ======================================================================
\section{Thảo luận}
\label{sec:discussion}

\tbd{Viết lại sau khi có N2--N5 và H1--H4.} Các điểm dự kiến thảo luận:
\begin{itemize}
    \item Nếu Bảng N2 vẫn cho thấy phụ thuộc kiến trúc sau khi thống nhất tầng chưng cất và chọn $\lambda_{sp}$ theo validation, thì kết luận của nghiên cứu trước chỉ đúng cho FB-tCNN và không nên khái quát cho CCKD nói chung.
    \item Quan hệ giữa $\Delta$CKA và mức cải thiện gợi ý một chỉ báo có thể đo trên validation trước khi chạy đầy đủ. Cần kiểm tra chỉ báo này có dự đoán được kết quả trên tập kiểm thử không.
    \item Kích thước tập kiểm thử của Benchmark (10 subject) là nguồn chính của sai lệch giữa một và năm cách chia.
\end{itemize}

\paragraph{Giới hạn.} Năm kiến trúc đều là mạng tích chập, chưa có mô hình dựa trên Transformer. Các tập điện cực thay thế được cố định trước và không bao phủ mọi tổ hợp. Hai bộ dữ liệu có thể có chung thiết bị thu và quy trình ghi \cantrich{kiểm tra thiết bị và quy trình ghi của Benchmark và BETA}, nên kết quả cần được kiểm tra thêm trên dữ liệu thu bằng thiết bị đeo. Các hạn chế kỹ thuật ở Mục~\ref{sec:setup} vẫn áp dụng.

% ======================================================================
\section{Kết luận}
\label{sec:conclusion}

\tbd{Viết lại sau khi có kết quả.} Bài báo đánh giá chưng cất quan hệ xuyên kênh cho SSVEP ít điện cực trên năm kiến trúc, hai bộ dữ liệu và năm cách chia subject. Với giao thức cũ, mức cải thiện phụ thuộc mạnh vào kiến trúc và bộ dữ liệu. Đánh giá trên một cách chia cho kết luận khác ở 5 trên 20 tổ hợp.

% ======================================================================
\bibliography{main}
\bibliographystyle{tmlr}

\appendix
\section{Phụ lục: số liệu theo từng cách chia}

\begin{table}[h]
\caption{Per-split CCKD gains for DeepConvNet, ShallowConvNet and Compact-CNN (earlier protocol, $\lambda_{sp}=100$, 8 seeds per split). $\bar\Delta \pm$ SD over five splits; $p_t$: one-sample $t$-test; bold if $p_t<0.05$.}
\label{tab:multisplit_gen}
\begin{center}
\footnotesize
\begin{tabular}{llccccccc}
\toprule
Backbone & Setting & 2026 & 3026 & 4026 & 5026 & 6026 & $\bar\Delta \pm$ SD & $p_t$ \\
\midrule
\multirow{4}{*}{DeepConvNet} & BM-3C & +0.52 & +0.45 & +0.53 & +0.67 & +0.79 & +0.59 $\pm$ 0.13 & \textbf{0.0006} \\
 & BM-1C & +1.17 & +0.45 & $-0.20$ & $-0.34$ & $-0.42$ & +0.13 $\pm$ 0.68 & 0.684 \\
 & BETA-3C & +0.90 & +1.16 & +0.92 & +0.96 & +0.28 & +0.84 $\pm$ 0.33 & \textbf{0.0047} \\
 & BETA-1C & +0.56 & +0.84 & +0.86 & +0.75 & +0.04 & +0.61 $\pm$ 0.34 & \textbf{0.0161} \\
\multirow{4}{*}{ShallowConvNet} & BM-3C & $-0.31$ & $-1.34$ & $-0.83$ & $-0.66$ & $-1.43$ & $-0.91 \pm 0.47$ & \textbf{0.0125} \\
 & BM-1C & $-0.55$ & $-0.78$ & $-0.43$ & $-1.30$ & $-1.69$ & $-0.95 \pm 0.53$ & \textbf{0.0161} \\
 & BETA-3C & +0.50 & +0.12 & +1.48 & +0.46 & $-0.27$ & +0.46 $\pm$ 0.65 & 0.191 \\
 & BETA-1C & +0.32 & $-0.14$ & +0.60 & +0.24 & $-0.16$ & +0.17 $\pm$ 0.32 & 0.300 \\
\multirow{4}{*}{Compact-CNN} & BM-3C & $-0.21$ & $-0.14$ & $-0.62$ & +1.36 & $-0.10$ & +0.06 $\pm$ 0.76 & 0.872 \\
 & BM-1C & $-0.52$ & $-0.47$ & $-0.20$ & $-0.43$ & +0.26 & $-0.27 \pm 0.32$ & 0.128 \\
 & BETA-3C & $-0.04$ & $-0.34$ & $-0.07$ & $-0.78$ & +0.04 & $-0.24 \pm 0.34$ & 0.190 \\
 & BETA-1C & $-0.12$ & $-0.63$ & $-0.22$ & $-0.27$ & +0.23 & $-0.20 \pm 0.31$ & 0.219 \\
\bottomrule
\end{tabular}
\end{center}
\end{table}

\begin{table}[h]
\caption{Per-split CCKD gains for EEGNet (earlier protocol, $\lambda_{sp}=100$, 8 seeds per split). Bold if $p_t<0.05$.}
\label{tab:eegnet_multisplit}
\begin{center}
\footnotesize
\begin{tabular}{lccccccc}
\toprule
Setting & 2026 & 3026 & 4026 & 5026 & 6026 & $\bar\Delta \pm$ SD & $p_t$ \\
\midrule
Benchmark 3C & +0.66 & +0.48 & +0.77 & +1.66 & +1.85 & \textbf{+1.09 $\pm$ 0.63} & \textbf{0.0179} \\
Benchmark 1C & $-0.17$ & $-1.83$ & $-0.96$ & $-0.31$ & $-0.03$ & $-0.66 \pm 0.75$ & 0.120 \\
BETA 3C & $-1.30$ & $-1.30$ & $-1.23$ & $-0.92$ & $-1.73$ & \textbf{$-1.30 \pm 0.29$} & \textbf{0.0006} \\
BETA 1C & $-1.02$ & $-1.52$ & $-0.60$ & $-1.55$ & $-1.75$ & \textbf{$-1.29 \pm 0.47$} & \textbf{0.0036} \\
\bottomrule
\end{tabular}
\end{center}
\end{table}

\begin{table}[h]
\caption{Sensitivity to $\lambda_{sp}$ (FB-tCNN, earlier protocol, 1.0 s, split 2026, 8 seeds). Val: mean validation accuracy over seeds (\%); Test: mean test accuracy (\%); $\Delta$: paired gain over the baseline of the same run (pp) with Wilcoxon $p$. $\star$: value selected on validation.}
\label{tab:lambda_val}
\begin{center}
\footnotesize
\begin{tabular}{llccccc}
\toprule
Dataset & Student & $\lambda_{sp}$ & Val (\%) & Test (\%) & $\Delta$ (pp) & $p$ \\
\midrule
\multirow{4}{*}{Benchmark} & \multirow{4}{*}{3C} & 1$^\star$ & 53.66 & 51.23 & +0.22 & 0.594 \\
 & & 10 & 53.42 & 51.55 & +0.51 & 0.055 \\
 & & 100 & 53.05 & 51.44 & +0.44 & 0.496 \\
 & & 1000 & 51.95 & 50.77 & $-0.27$ & 0.820 \\
\midrule
\multirow{4}{*}{Benchmark} & \multirow{4}{*}{1C} & 1$^\star$ & 49.45 & 44.35 & $-0.33$ & 0.090 \\
 & & 10 & 49.36 & 44.68 & +0.18 & 0.652 \\
 & & 100 & 49.41 & 45.35 & +0.73 & 0.027 \\
 & & 1000 & 48.27 & 43.97 & $-0.68$ & 0.322 \\
\midrule
\multirow{4}{*}{BETA} & \multirow{4}{*}{3C} & 1 & 46.19 & 38.26 & $-0.18$ & 0.927 \\
 & & 10 & 46.23 & 38.40 & $-0.03$ & 0.985 \\
 & & 100$^\star$ & 47.97 & 39.32 & +0.93 & 0.026 \\
 & & 1000 & 47.52 & 38.67 & +0.25 & 0.681 \\
\midrule
\multirow{4}{*}{BETA} & \multirow{4}{*}{1C} & 1 & 41.27 & 35.14 & +0.03 & 0.931 \\
 & & 10 & 41.69 & 34.98 & $-0.23$ & 0.538 \\
 & & 100$^\star$ & 42.61 & 35.54 & +0.40 & 0.173 \\
 & & 1000 & 40.25 & 33.54 & $-1.53$ & 0.017 \\
\bottomrule
\end{tabular}
\end{center}
\end{table}

\end{document}
